% Формат файла
\documentclass[12pt]{article}
\usepackage[paperheight=297mm,
	paperwidth=210mm,
	top=20mm,
	bottom=20mm,
	left=15mm,
right=15mm]{geometry}

% Импорт преамбулы
\usepackage{../mypreamble}

% Оглавление
\title{\textbf{Дискретная математика}}\date{}\author{}

\begin{document}

\maketitle
\tableofcontents
\label{toc}
\newpage

\section{Бинарные отношения}

\begin{definition}
	Бинарным отношением на множествах $A, B$ называется $R\subseteq
	A\times B$. Если $(x,y)\in R$, говорят, что $x$ и $y$ находятся в
	отношении $R$.
\end{definition}

\subsection{Функции}

\begin{definition}
	Функцией из множества $A$ в множество $B$ называется такое бинарное
	отношение $f \subseteq A\times B$, что для каждого $x\in A$
	существует не более одной пары $(x, y)\in f$. Элементы множества $A$
	называют аргументами функции, а $B$ -- её значениями.
\end{definition}
\begin{definition}
	Пусть $f\colon A\to B$. Областью определения $f$ называется:
	$$\dom(f)\coloneqq\{x\in A\mid \exists\, y\in B\colon y=f(x)\}$$
\end{definition}
\begin{definition}
	Функция $f\colon A\to B$ называется тотальной, если $\dom(f)=A$.
\end{definition}
\begin{definition}
	Отображением называется тотальная функция.
\end{definition}
\begin{definition}
	Пусть $f\colon A\to B$. Областью значений $f$ называется:
	$$\range(f)\coloneqq \{y\in B\mid \exists\, x\in A\colon y=f(x)\}$$
\end{definition}
\begin{definition}
	Булеаном (или степенью множества) $A$ называется множество
	$\mathcal{P}(A)$ всех подмножеств множества $A$, при этом
	$\varnothing, A\in \mathcal{P}(A)$.
\end{definition}

\subsection{Последовательности}

\begin{definition}
	Начальным отрезком натурального ряда называется:
	$$[n]\coloneqq \{x\colon x<n, x\in \N\}$$
\end{definition}
\begin{definition}
	Конечной последовательностью элементов множества $A$ (или словом в
	алфавите $A$) называется тотальная функция $[n]\to A$. Тогда $n$
	называется длиной последовательности.
\end{definition}
\begin{definition}
	Бесконечной последовательностью элементов множества $A$ (или
	сверхсловом в алфавите $A$) называется тотальная функция $\N\to A$.
\end{definition}

\subsection{Композиции функций}

\begin{definition}
	Композицией функций $f\colon A\to B$ и $g\colon B\to C$
	называется $g\circ f\colon A\to C$, которая определена на
	$x\in \dom f\colon f(x) \in \dom g$ и равна $g(f(x))$.
\end{definition}
\begin{definition}
	Для множества $A$ тождественным отображением называется $\id_A\colon
	x\mapsto x$, переводящие каждый элемент множества само в себя.
\end{definition}
\begin{definition}
	Для тотальной функции $f\colon A\to B$ левой обратной называется
	такая тотальная функция $g\colon B\to A$, что $g\circ f=\id_A$.
\end{definition}
\begin{definition}
	Для тотальной функции $f\colon A\to B$ правой обратной называется
	такая тотальная функция $g\colon B\to A$, что $f\circ g=\id_B$.
\end{definition}
\begin{definition}
	Тотальная функция $f\colon A\to B$ называется инъекцией, если:
	$$\forall x_1, x_2 \in A\ x_1 \neq x_2\implies f(x_1)\neq f(x_2)$$
	Равносильное определение получается контрапозицией:
	$$\forall x_1, x_2 \in A\ f(x_1)=f(x_2)\implies x_1=x_2$$
\end{definition}
\begin{lemma}
	У тотальной функции $f\colon A\to B$ есть левая обратная $\iff$ $f$
	-- инъекция.
	\begin{proof}
		Пусть $g \circ f=\id_A$. Тогда $f(x_1)=f(x_2)\implies
		x_1=g(f(x_1))=g(f(x_2))=x_2$, то есть $f$ -- инъекция.\bigskip

		Пусть $f$ -- инъекция. Рассмотрим такую $g\colon B\to A$, что
		$g(y)=x\colon f(x)=y$, если $y\in \range f$, и $g(y)$ произвольное,
		если $y\notin \range f$. Тогда $\forall x \in A\ g(f(x))=x$, то
		есть $g$ -- левая обратная.
	\end{proof}
\end{lemma}
\begin{definition}
	Тотальная функция $f\colon A\to B$ называется сюръекцией, если $\range f=B$.
\end{definition}
\begin{lemma}
	У тотальной функции $f\colon A\to B$ есть правая обратная $\iff$ $f$
	-- сюръекция.
	\begin{proof}
		Пусть $g\colon B\to A$ -- правая обратная. Тогда $\forall y\in
		B\ f(g(y))=y$, то есть:
		$$\forall y \in B\ \exists\, x\in A\colon f(x)=y$$
		Пусть $f$ -- сюръекция. Рассмотрим такую $g\colon B\to A$, что
		$g(y)=x\colon f(x)=y$. Если таких $x$ несколько, то берём
		произвольный. Тогда $\forall y\in B\ f(g(y))=y$, то есть $g$ --
		правая обратная.
	\end{proof}
\end{lemma}
\begin{definition}
	Тотальная функция называется биекцией (или взаимно-однозначным
	соответствием), если она инъекция и сюръекция.
\end{definition}
\begin{definition}
	Обратной функцией (или обратным отображением) к биекции $f\colon
	A\to B$ называется такая $f^{-1}\colon B\to A$, что если $f\colon
	x\mapsto y$, то $f^{-1}\colon y \mapsto x$.
\end{definition}
\begin{lemma}
	Пусть $f\colon A\to B$, $g\colon B\to C$. Если $f, g$ -- инъекции,
	то $g\circ f$ -- инъекция. Если $f, g$ -- сюръекции, то $g\circ f$
	-- срюръекция. Если $f, g$ -- биекции, то $g\circ f$ -- биекция.
	\begin{proof}
		Пусть $(g\circ f)(x_1)=(g\circ f)(x_2)$, то есть
		$g(f(x_1))=g(f(x_2))$. В силу инъективности $g$ $f(x_1)=f(x_2)$, но
		тогда $x_1=x_2$, так как $f$ -- инъекция. Отсюда $g\circ f$ --
		инъекция.\bigskip

		$g$ -- сюръекция, значит $\forall z\in C\ \exists\,y\in B\colon
		g(y)=z$. Так как $f$ -- сюръекция, $\forall y\in B\ \exists\,x\in
		A\colon f(x)=y$. То есть $\forall z\in C\ \exists\,x\in A\colon
		g(f(x))=z$, значит $g\circ f$ -- сюръекция.\bigskip

		Если $f, g$ -- инъекции и сюръекции, то из первых двух
		доказательств $g\circ f$ также инъекция и сюръекция, то есть биекция.
	\end{proof}
\end{lemma}

\section{Принцип математической индукции}

\begin{axiom}[Принцип математической индукции]
	Пусть для последовательности утверждений $A_1,\ldots,A_n$ известно,
	что $A_1$ истинно и $\forall n\ A_n\to A_{n+1}$. Тогда $\forall n\
	A_n$ истинно. $A_1$ называют базой индукции. Посылку импликации
	$A_n\to A_{n+1}$ называют индуктивным предположением.
\end{axiom}
\begin{statement}[Неравенство Бернулли]
	$$(1+x)^n\geq 1+nx\ \forall n\geq 1, x\geq -1$$
	\begin{proof}
		Заметим, что для $n=1$ верно:
		$$1+x\geq 1 + x$$
		Пусть для $n$ верно, рассмотрим для $n+1$:
		$$(1+x)^{n+1}=(1+x)^n(1+x)\geq (1+nx)(1+x)=1+(n+1)x+nx^2\geq 1+(n+1)x$$
	\end{proof}
\end{statement}
\begin{theorem}[Формула для биномиальных коэффициентов]
	$$\binom{n}{k}=\frac{n!}{k!(n-k)!}, 0\leq k \leq n$$
	\begin{proof}
		Заметим, что для $n=1$ верно:
		$$(1+x)^1=\binom{1}{0}+\binom{1}{1}x=1+x$$
		Пусть для $n$ верно, рассмотрим для $n+1$. Заметим, что для $k=0$ и
		$k=n+1$ верно:
		$$\binom{n+1}{0}=1=\frac{(n+1)!}{0!(n+1)!}\qquad
		\binom{n+1}{n+1}=1=\frac{(n+1)!}{(n+1)!0!}$$
		Сравним коэффициенты при $x^{n+1-k}y^k$ в левой и правой частях равенства:
		$$(1+x)^{n+1}=(1+x)^n(1+x)$$
		Получаем для $0<k<n+1$ равенства:
		\begin{equation}
			\tag{2.2}
			\label{Сумма биномиальных коэффициентов}
			\binom{n+1}{k}=\binom{n}{k}+\binom{n}{k-1}
		\end{equation}
		Таким образом, применив формулу для $n$, получаем, что:
		$$\frac{n!}{k!(n-k)!}+\frac{n!}{(k-1)!(n-k+1)!}=\frac{n!(n-k+1+k)}{k!(n-k+1)!}=\frac{(n+1)!}{k!(n+1-k)!}$$
	\end{proof}
\end{theorem}
\begin{theorem}[Принцип полной математической индукции]
	Пусть для последовательности утверждений $A_0,A_1,\ldots,A_n$ верно, что
	$\forall i<n\ A_i\to A_n$. Тогда $\forall n\ A_n$ истинно.
	\begin{proof}
		Пусть $B_n=A_0\land A_1\land \ldots \land A_n$. Индуктивное
		предположение гласит, что $B_0=A_0$ истинно и $\forall n\ B_n\to
		A_{n+1}$, но $B_n\to B_{n+1}$, так как $B\to A \equiv B\to B\land
		A$. Тогда по принципу математической индукции $\forall n\ B_n$
		истинно, а значит и $A_n$ истинно.
	\end{proof}
\end{theorem}
\begin{theorem}[Формула Бине]
	$$F_n=\frac{\varphi^{n+1}-\psi^{n+1}}{\sqrt{5}}\text{, где
	}\varphi=\frac{1+\sqrt{5}}{2}, \psi=\frac{1-\sqrt{5}}{2}$$
	\begin{proof}
		Покажем, что для $n=0$ и $n=1$ верно:
		$$F_0=\frac{\varphi-\psi}{\sqrt{5}}=\frac{\sqrt{5}}{\sqrt{5}}=1$$
		$$F_1=\frac{\varphi^2-\psi^2}{\sqrt{5}}=\frac{(\varphi-\psi)(\varphi+\psi)}{\sqrt{5}}=\frac{\sqrt{5}}{\sqrt{5}}=1$$
		Пусть формула верна для $n$ и $n+1$, покажем, что она верна для $n+2$:
		$$F_{n+2}=F_{n}+F_{n+1}=\frac{\varphi^{n+1}-\psi^{n+1}}{\sqrt{5}}+\frac{\varphi^{n+2}-\psi^{n+2}}{\sqrt{5}}=\frac{1}{\sqrt{5}}\left(\varphi^{n+1}(\varphi+1)-\psi^{n+1}(\psi+1)\right)$$
		Заметим, что $\varphi+\psi=1, \varphi\cdot\psi=-1$, то есть они
		являются корнями уравнения $x^2-x-1=0$, значит $\varphi^2=\varphi+1,
		\psi^2=\psi+1$. Применим это к формуле для $F_{n+2}$:
		$$\frac{1}{\sqrt{5}}\left(\varphi^{n+1}(\varphi+1)-\psi^{n+1}(\psi+1)\right)=\frac{\varphi^{n+3}-\psi^{n+3}}{\sqrt{5}}$$
	\end{proof}
\end{theorem}
\begin{theorem}[Принцип наименьшего числа]
	Любое непустое подмножество натуральных чисел содержит минимальный элемент.
	\begin{proof}
		Рассмотрим контрапозицию: $\forall a\in X\ \exists\,b\in X\colon
		b<a\implies X=\varnothing$, то есть $\forall n\in \N\ n\notin X$.
		Тогда заметим, что верно $\forall i<n\ i\notin X\to n\notin X$,
		иначе $n$ -- минимальный элемент.
	\end{proof}
\end{theorem}
\begin{statement}
	Соседние числа Фибоначчи взаимно простые.
	\begin{proof}
		Рассмотрим $B=\{n\colon \exists\, d\colon d\mid F_n, d\mid F_{n+1},
		d>1\}$, тогда $k\neq0$ (так как $F_0=1$) -- наименьший элемент
		этого множества. $F_{k-1}=F_{k+1}-F_k$ должно делиться на $d$, то
		есть $k-1\in B$, но $k-1<k$ и $k$ -- наименьший, противоречие.
	\end{proof}
\end{statement}

\section{Перечислительная комбинаторика}

\begin{theorem}[Правило суммы]
	Для конечных непересекающихся множеств $A, B$ выполняется:
	$$A\cap B=\varnothing\implies |A\cup B|=|A|+|B|$$
\end{theorem}
\begin{definition}[Декартово произведение]
	$$A_1\times A_2\times \ldots \times A_n=\{(a_1, a_2, \ldots,
	a_n)\mid a_i \in A_i\}$$
\end{definition}
\begin{lemma}[Правило произведения]
	$$|A_1\times A_2\times \ldots \times
	A_n|=|A_1|\cdot|A_2|\cdot\ldots\cdot |A_n|$$
	\begin{proof}
		Докажем по индукции. Заметим, что для $n=1$ верно: $|A_1|=|A_1|$.
		Пусть верно для $n$, рассмотрим для $n+1$:
		\begin{enumerate}[itemsep=0pt]
			\item Если $A_{n+1}=\varnothing$, то выполняется:
				$$|A_1\times A_2\times \ldots \times
				A_{n+1}|=|\varnothing|=0=|A_1|\cdot|A_2|\cdot\ldots\cdot |A_n|\cdot 0$$
			\item Если $A_{n+1}=\{a\}$, то каждая последовательность для $n$
				однозначно дополняется до $n+1$ дописыванием последнего члена
				$a$. Также в обратную сторону для любой последовательности длины
				$n+1$ можно получить последовательность длины $n$, вычеркнув
				последний член последовательности. То есть:
				$$|A_1\times A_2\times \ldots \times
				A_{n+1}|=|A_1\times A_2\times \ldots \times
				A_{n}|=|A_1|\cdot|A_2|\cdot\ldots\cdot |A_n|\cdot
				1=|A_1|\cdot|A_2|\cdot\ldots\cdot |A_n|\cdot
				|A_{n+1}|$$
			\item Если $A_{n+1}=\{u_1,\ldots,u_l\}, |A_{n+1}|=k\geq 2$, то
				разложим декартово произведение в объединение непересекающихся множеств:
				$$A_1\times A_2\times \ldots \times A_{n+1}=\bigcup_{i=1}^k A_i',\  A_i'=
				A_1\times A_2\times \ldots \times A_{n}\times \{u_i\}$$
				Из предыдущего пункта нам известно, что:
				$$|A_i'|=|A_1\times A_2\times \ldots \times A_{n}|$$
				Тогда по правилу суммы получаем, что:
				$$|A_1\times A_2\times \ldots \times
				A_{n+1}|=\sum_{i=1}^{k}|A_1\times A_2\times \ldots \times
				A_{n}|=|A_1|\cdot|A_2|\cdot\ldots\cdot |A_n|\cdot
				k=|A_1|\cdot|A_2|\cdot\ldots\cdot |A_n|\cdot
				|A_{n+1}|$$
		\end{enumerate}
	\end{proof}
\end{lemma}
\begin{definition}
	Слово в алфавите $A$ -- конечная последовательность элементов $A$.
	$A^*$ -- множество всех слов в алфавите $A$. $A^n$ -- множество всех
	слов длины $n$ в алфавите $A$.
\end{definition}
\begin{theorem}\label{Количество слов в алфавите}
	Количество слов длины $n$ в алфавите из $k$ букв равно $k^n$.
	\begin{proof}
		Пусть $A$ -- алфавит. Тогда $A^n=A\times A\times \ldots \times A$
		(умножение на $A$ $n$ раз), но по правилу произведения $|A^n|=|A|^n=k^n$.
	\end{proof}
\end{theorem}
\begin{statement}\label{Прообраз непересекающихся не пересекается}
	Пусть $f\colon A\to B, Y_1\subseteq B, Y_2\subseteq B, Y_1\cap
	Y_2=\varnothing$. Тогда $f^{-1}(Y_1)\cap f^{-1}(Y_2)=\varnothing$.
	\begin{proof}
		У каждого аргумента функции есть не более одного значения, поэтому
		если $f(a)\in Y_1$, то $a\notin f^{-1}(Y_2)$ и наоборот.
	\end{proof}
\end{statement}
\begin{statement}\label{Мощность полного прообраза равна мощности
	суммы прообразов}
	Пусть $f\colon A\to B, Y\subseteq B$, тогда:
	$$|f^{-1}(Y)|=\sum_{b\in Y}|f^{-1}(\{b\})|$$
	\begin{proof}
		Из утверждения \ref{Прообраз непересекающихся не пересекается}
		прообразы элементов из $Y$ не пересекаются. Тогда можно применить
		правило суммы.
	\end{proof}
\end{statement}
\begin{lemma}
	Пусть $f\colon A\to B$, $A$ и $B$ конечные, тогда выполняется:
	\begin{center}
		\begin{minipage}{0.45\textwidth}
			\begin{enumerate}[itemsep=0pt]
				\item Если $f$ -- инъекция, то $|A|\leq |B|$;\label{Зависимость
					количества элементов от вида функции}
				\item Если $f$ -- сюръекция, то $|A|\geq |B|$;
				\item Если $f$ -- биекция, то $|A|= |B|$.
			\end{enumerate}
		\end{minipage}
	\end{center}
	\begin{proof}
		Пусть $a_i=|f^{-1}(\{i\})|, i\in B$. Если $f$ -- инъекция, то
		$a_i\leq1\ \forall i\in B$. Тогда из утверждения \ref{Мощность
		полного прообраза равна мощности суммы прообразов} получаем:
		$$|A|=\sum_{i\in B}a_i\leq \sum_{i\in B}1=|B|$$
		Аналогично если $f$ -- сюръекция, то $a_i\geq1\ \forall i\in B$,
		значит $|A|\geq |B|$. Если $f$ -- биекция, то $a_i=1\ \forall i\in
		B$, то есть $|A|=|B|$.
	\end{proof}
\end{lemma}
\setcounter{subsection}{6}
\begin{consequence}[Принцип Дирихле]
	Если $n$ кроликов посажены в $m<n$ клеток, то хотя бы два кролика
	сидят в одной клетке.
	\begin{proof}
		Рассадка кроликов -- это тотальная функция $f\colon
		\text{кролики}\to\text{клетки}$. Рассмотрим контрапозицию к пункту
		\ref{Зависимость количества элементов от вида функции} предыдущей
		леммы: Если $|A|>|B|$, то $f$ -- не инъекция. То есть для двух
		аргументов (кроликов) значение функции (клетка) совпадают.
	\end{proof}
\end{consequence}
\begin{consequence}\label{Счет гвоздей в упаковках}
	Пусть $f\colon A\to B$ -- такая сюръекция, что:
	$$\forall b_1, b_2\in
	B\ |f^{-1}(\{b_1\})|=|f^{-1}(\{b_2\})|=k\text{, тогда }|A|=k|B|$$
\end{consequence}
\setcounter{subsection}{0}
\begin{statement}[Количество неупорядоченных пар]
	$$C_n^2=\{\{x,y\}\mid x, y\in [n], x\neq y\},\quad|C_n^2|=\frac{n(n-1)}{2}$$
	\begin{proof}
		Пусть $A_n^2=\{(x,y)\mid x, y\in [n], x\neq y\}$. По правилу
		произведения всего упорядоченных пар $n^2$, из них $n$ состоит из
		одинаковых элементов, то есть $|A_n^2|=n^2-n=n(n-1)$. Определим
		$f\colon A_n^2\to C_n^2$, которая «забывает» порядок в паре, тогда:
		$$\forall \{x, y\}\in C_n^2\ f^{-1}(\{x, y\})=\{(x,y), (y, x)\}$$
		Отсюда из следствия \ref{Счет гвоздей в упаковках}
		$|A_n^2|=2|C_n^2|$, то есть:
		$$|C_n^2|=\frac{n(n-1)}{2}$$
	\end{proof}
\end{statement}
\begin{lemma}\label{Количество тотальных функций}
	$B^A$ -- множество тотальных функций $A\to B$. Тогда для конечных
	множеств: $$|B^A|=|B|^{|A|}$$
	\begin{proof}
		Пусть $|A|=n, |B|=k$. Заметим, что слово в алфавите $B$ длины $n$
		-- это последовательность элементов $B$, то есть тотальная функция
		$[n] \to B$. Количество таких слов по теореме \ref{Количество слов
		в алфавите} равно $k^n$. Докажем, что количество таких функций
		совпадает с количеством тотальных функций $A\to B$.\bigskip

		Поскольку $|A|=n$, существует биекция $\alpha\colon [n]\to A$.
		Пусть $a_i=\alpha(i)$, то есть $A=\{a_0,\ldots,a_{n-1}\}$. Теперь
		сопоставим тотальной функции $f\colon A\to B$ композицию
		$\beta(f)=f\circ\alpha$. Это функция вида $[n]\to B$, то есть слово
		$b_0b_1\ldots b_{n-1}$ в алфавите $B$, где $b_i=f(a_i)$.\bigskip

		Тогда разным функциям $f, g$ сопоставлены разные слова: $f(a_j)\neq
		g(a_j)\implies \beta(f)_j\neq \beta(g)_j$, то есть $\beta$ --
		инъекция. При этом по каждому слову $b_0b_1\ldots b_{n-1}$
		однозначно определяется такая $f$, что $\beta(f)=b$, это $f\colon
		f(a_j)=b_j$, то есть $\beta$ -- сюръекция, а значит и биекция, то
		есть $|B^A|$ совпадает с количеством слов длины $n$ в алфавите из
		$k$ элементов, значит $|B^A|=k^n=|B|^{|A|}$.
	\end{proof}
\end{lemma}
\begin{statement}
	Количество функций $A\to B$, $|A|=n, |B|=k$, равно $(k+1)^n$.
	\begin{proof}
		Доказательство аналогично доказательству леммы \ref{Количество
		тотальных функций}, но теперь для биекции рассматриваем
		$B'=B\cup\{\texttt{void}\}$, где элементу \texttt{void} будут
		соответствовать все $A\to B$, которые не определены.
	\end{proof}
\end{statement}

\subsection{Размещения}

\begin{definition}
	Размещением $A_n^k$ называется упорядоченная выборка или слово длины
	$k$ в алфавите из $n$ элементов, где все символы разные. Если не
	оговорено иного, алфавитное множество имеет вид $\{1, 2,\ldots, n\}$.
\end{definition}
\begin{statement}\label{Обрезка слов}
	$$A_n^k=(n-k+1)A_n^{k-1}$$
	\begin{proof}
		Рассмотрим $\tau\colon A_n^k\to A_n^{k-1}$ -- функция обрезки
		последнего символа в слове. Тогда $|\tau^{-1}(x)|=n-k+1$, так как
		обрезать можно любой символ их тех, что ещё не встречались в слове.
	\end{proof}
\end{statement}
\begin{theorem}
	$$A_n^k=\frac{n!}{(n-k)!}$$
	\begin{proof}
		Из утверждения \ref{Обрезка слов} получаем:
		$$A_n^k=(n-k+1)A_n^{k-1}=(n-k+1)(n-k+2)A_n^{k-2}=\ldots=(n-k+1)\cdot\ldots\cdot
		n\cdot A_n^0$$
		При этом $A_n^0=1$, так как пустое слово есть только одно. Отсюда
		получаем искомую формулу.
	\end{proof}
\end{theorem}
\begin{statement}
	Количество инъекций $A\to B$, $|A|=n, |B|=k$, равно $A_n^k$.
	\begin{proof}
		Количество таких инъекций равно количеству инъекций $[k]\to[n]$, то
		есть количеству слов длины $k$ из алфавита длины $n$, где все
		символы разные, а оно равно $A_n^k$ по определению.
	\end{proof}
\end{statement}
\setcounter{subsection}{12}
\begin{consequence}
	Количество биекций на $n$-элементном множестве равно $n!$.
	\begin{proof}
		Заметим, что если $f\colon [n]\to[n]$ -- инъекция, то $f$ также
		сюръекция, а значит количество биекций равно $A_n^n=n!$.
	\end{proof}
\end{consequence}
\setcounter{subsection}{1}
\begin{definition}
	Размещение из $n$ по $n$ называется перестановкой.
\end{definition}
\begin{definition}
	Индикаторной (или характеристической) функцией называется:
	$$\chi_S(x)=
	\begin{+cases}
		1, x\in S\\
		0, x\notin S
	\end{+cases}$$
\end{definition}
\begin{statement}\label{Биекции между словами, тотальными функциями и
	подмножествами}
	Пусть $|X|=n$, тогда $|\mathcal{P}(X)|=2^n$.
	\begin{proof}
		Пусть $S\subseteq X$. Заметим, что $\chi_A=\chi_B\iff A=B$, так как
		по определению $\chi_A=\chi_B$ значит, что каждый элемент либо
		входит в $A\cap B$, либо, не входит ни в $A$, ни в $B$. Тогда
		$S\mapsto \chi_S$ -- инъекция, докажем, что это сюръекция.\bigskip

		Рассмотрим тотальную $f\colon X\to \{0,1\}$. Пусть $S_f=\{x\mid
		f(x)=1\}$. Тогда $f=\chi_{S_f}$, то есть $S\mapsto \chi_S$ --
		сюръекция, а значит и биекция.\bigskip

		Таким образом, $S\mapsto \chi_S$ -- биекция $\mathcal{P}(X)\to \{0,
		1\}^X$, тогда по лемме \ref{Количество тотальных функций}
		$|\mathcal{P}(X)|=2^n$.
	\end{proof}
\end{statement}

\subsection{Сочетания}

\begin{definition}
	Сочетанием $C_n^k$ называется неупорядоченная выборка или
	подмножество $n$-элементного множества, в котором $k$ элементов.
	Если не оговорено иного, $n$-элементное множество имеет вид $\{1,
	2,\ldots, n\}$.
\end{definition}
\begin{statement}\label{Количество двоичных слов с k единицами}
	Количество двоичных слов длины $n$ с $k$ единицами равно $C_n^k$.
	\begin{proof}
		В доказательстве утверждения \ref{Биекции между словами, тотальными
		функциями и подмножествами} мы установили биекции между двоичными
		словами длины $n$, тотальными функциями $[n]\to\{0,1\}$ и
		подмножествами $n$-элементного множества. Отсюда двоичному слову с
		$k$ единицами соответствует индикаторная функция, в значениях
		которой $k$ единиц, а ей соответствует подмножество размера $k$, то
		есть число сочетаний совпадает с числом двоичных слов из условия.
	\end{proof}
\end{statement}
\begin{theorem}
	$$C_n^k=\frac{n!}{k!(n-k)!}$$
	\begin{proof}
		Рассмотрим $f\colon A_n^K\to C_n^k$ -- функцию, которая «забывает»
		порядок. Эта функция сюръективна, при этом $|f^{-1}(\{S\})|$, где
		$S\in C_n^k$, равно количеству способов упорядочить $k$ элементов,
		то есть количеству перестановок $A_k^k=k!$. Тогда из следствия
		\ref{Счет гвоздей в упаковках} получаем:
		$$C_n^k=k!\cdot A_n^k=\frac{n!}{k!(n-k)!}$$
	\end{proof}
\end{theorem}
\setcounter{subsection}{14}
\begin{consequence}
	$$\binom{n}{k}=C_n^k$$
	\textit{Замечание:} Равенство также можно доказать, представив
	биномиальный коэффициент в разложении $(x+y)^n$ как количество слов
	алфавита $\{x,y\}$, в которых количество символов $x$ равно $k$, а
	символов $y$ равно $n-k$ (так как других символов там нет).
\end{consequence}
\setcounter{subsection}{2}

\subsection{Биномиальные коэффициенты}

\begin{statement}
	Любое число в треугольнике Паскаля кроме крайних единиц равно сумме
	соседних чисел, которые стоят над ним.
	\begin{proof}
		Следует из равенства \ref{Сумма биномиальных коэффициентов}.
	\end{proof}
\end{statement}
\begin{statement}
	Пусть $T(a, b)$ -- количество монотонных путей из начала координат в
	$(a, b)$, тогда:
	$$T(a, b)=\binom{a+b}{a}=\frac{(a+b)!}{a!\cdot b!}$$
	\begin{proof}
		Построим биекцию между монотонными путями и двоичными словами в
		алфавите $\{x, y\}$ длины $a+b$ с ровно $a$ символами $x$. Тогда
		символ $x$ в слове будет означать движение вправо, а $y$ -- вверх.
		Для каждого слова можно однозначно задать монотонный путь и
		наоборот, значит это действительно биекция, а из утверждения
		\ref{Количество двоичных слов с k единицами} количество таких слов
		равно $C_{a+b}^a$, то есть биномиальному коэффициенту из условия.
	\end{proof}
\end{statement}
\begin{statement}
	$$\binom{n}{k}=\binom{n}{n-k}$$
	\begin{proof}[Доказательство по формуле]
		$$\binom{n}{k}=\frac{n!}{k!(n-k)!}=\binom{n}{n-k}$$
	\end{proof}
	\begin{proof}[Доказательство через Бином Ньютона]
		Рассмотрим $(x+y)^n$ и $(y+x)^n$. Поскольку они равны, должно выполняться:
		$$\binom{n}{k}x^{n-k}y^k=\binom{n}{n-k}x^{n-k}y^k$$
	\end{proof}
	\begin{proof}[Комбинаторное доказательство]
		Рассмотрим биекцию $k$-элементного подмножества\\$n$-элементного
		множества с его дополнением, в котором
		$n-k$ элементов. Тогда $C_n^k=C_n^{n-k}$, а значит равны и
		биномиальные коэффициенты.
	\end{proof}
\end{statement}
\setcounter{subsection}{18}
\begin{consequence}
	Строки треугольника Паскаля симметричны относительно середины.
\end{consequence}
\setcounter{subsection}{3}
\begin{statement}
	Сумма чисел в $n$-ой строке треугольника Паскаля равна $2^n$.
	\begin{proof}[Аналитическое доказательство]
		Запишем формулу бинома Ньютона. При $x=y=1$ получаем сумму
		биномиальных коэффициентов:
		$$(x+y)^n=\sum_{k=0}^{n}\binom{n}{k}x^{n-k}y^k\qquad2^n=\sum_{k=0}^{n}\binom{n}{k}$$
	\end{proof}
	\begin{proof}[Комбинаторное доказательство]
		Разобьём двоичные слова на группы по количеству единиц в них, то
		есть в $k$-ю группу попадаются слова с $k$ единицами. В каждой
		группе $C_n^k$ слов, а всего слов в группах, то есть двоичных слов
		длины $n$ ровно $2^n$.
	\end{proof}
\end{statement}
\begin{statement}
	При $n>0$ количество подмножеств $n$-элементного множества с чётным
	числом элементов равно количеству подмножеств с нечётным числом элементов.
	\begin{proof}[Аналитическое доказательство]
		$$(1-1)^n=\sum_{k=0}^{n}\binom{n}{k}(-1)^k=\sum_{k\text{
		чётное}}\binom{n}{k}-\sum_{k\text{ нечётное}}\binom{n}{k}=0$$
	\end{proof}
	\begin{proof}[Комбинаторное доказательство]
		Разобьём подмножества множества $[n]$ на пары вида:
		$$\{\{n-1\}\cup S, S\},\ S\subseteq [n-1]$$
		Каждое множество входит ровно в одну пару. При этом количество
		элементов в множествах в паре отличаются ровно на 1, то есть одно
		из них содержит чётное количество элементов, а другое -- нечётное.
		Таким образом, мы сопоставили все множества с чётным
		числом элементов множествам с нечётным числом элементов, а значит
		их равное количество.
	\end{proof}
\end{statement}
\begin{statement}\label{Числа в треугольнике Паскаля возрастают к середине}
	В первой половине строки треугольника Паскаля числа возрастают.
	\begin{proof}[Аналитическое доказательство]
		Запишем условие возрастания значений в строке:
		$$1<\left.\binom{n}{k+1}\middle/\binom{n}{k}\right.=\frac{n!}{(k+1)!(n-k-1)!}\cdot
		\frac{k!(n-k)!}{n!}=\frac{n-k}{k+1}\iff 2k+1<n$$
	\end{proof}
	\begin{proof}[Комбинаторное доказательство]
		Будем рассматривать количество двоичных слов длины $n$, содержащих
		ровно $k$ единиц (имеющих вес $k$). Пусть $u_0u_1\ldots u_{n-1}$ --
		слово веса $k$. Введём $Z_p$ -- разность количества нулей и
		количества единиц в префиксе $u_0u_1\ldots u_p$ данного слова. В
		силу того, что мы рассматриваем первую половину строки, $2k<n-1$,
		значит $k<n-k$, то есть в слове нулей больше, чем единиц. Тогда
		$Z_{n-1}\geq 1$.\bigskip

		Пусть $q$ -- самый маленький индекс, при котором $Z_q=1$.
		Сопоставим каждому слову $u=u_0u_1\ldots u_{n-1}$ новое слово
		$v=\bar{u}_0\bar{u}_1\ldots \bar{u}_qu_{q+1}\ldots u_{n-1}$,
		полученное инвертированием всех символов до $q$. Мы выбирали $q$
		таким образом, что в первых $q$ символах слова нулей на 1 больше,
		чем единиц. Тогда после инвертирования уже единиц будет на 1
		больше. То есть мы построили функцию $\alpha\colon u\mapsto v$,
		которая переводит слова веса $k$ в слова веса $k+1$. Покажем, что
		это инъекция.\bigskip

		Заметим, что $Z_i(u)=-Z_i(v), 0\leq i \leq q$. Тогда если $v\in
		\range \alpha$, можно однозначно определить $u$, инвертировав
		первые $q'$ символов, где $q'$ -- минимальный индекс, при котором
		$Z_{q'}(v)=-1$. Если такой индекс не встречается, то $v\notin \range \alpha$.
	\end{proof}
\end{statement}
\begin{statement}
	$$\binom{2n}{n}\geq \frac{2^{2n}}{2n+1}$$
	\begin{proof}
		Сумма биномиальных коэффициентов равна $2^{2n}$. Самый большой из
		них -- средний, а всего коэффициентов $2n+1$, то есть:
		$$(2n+1)\binom{2n}{n}\geq 2^{2n}$$
	\end{proof}
\end{statement}

\subsection{Числа Каталана}

\begin{definition}
	Число Каталана $C_n$ -- это количество последовательностей Дика, то
	есть последовательностей из $n$ штук $+1$ и $n$ штук $-1$, в которых
	сумма по любому префиксу неотрицательна.
\end{definition}
\begin{theorem}
	$$C_n=\frac{1}{n+1}\binom{2n}{n}$$
	\begin{proof}
		Аналогично комбинаторному доказательству утверждения \ref{Числа в
		треугольнике Паскаля возрастают к середине} рассмотрим
		последовательности, но вместо нулей будут $+1$, а вместо единиц
		-- $-1$. В утверждении мы определили $\alpha$ таким образом, чтобы
		получать слова, в которых единиц больше нулей. То есть в новых
		обозначениях тех последовательностей, у которых сумма отрицательна.
		Значит нам нужны все те последовательности, которые не принадлежат
		$\range \alpha$:
		$$C_n=\binom{2n}{n}-\binom{2n}{n-1}=\frac{(2n)!}{n!\cdot
		n!}-\frac{(2n)!}{(n-1)!(n+1)!}=\frac{(2n)!}{n!\cdot
		n!}\left(1-\frac{n}{n+1}\right)=\binom{2n}{n}\frac{1}{n+1}$$
	\end{proof}
\end{theorem}
\begin{theorem}
	$$C_n=\sum_{k=0}^{n-1}C_k\cdot C_{n-k-1},\ C_0=1$$
	\begin{proof}
		Заметим, что последовательность Дика $x$ должна начинаться с нуля.
		Выберем её минимальный префикс, сумма элементов которого равна
		нулю. Такой префикс заканчивается на $-1$, ведь иначе сумма
		префикса до этого символа была отрицательной.\bigskip

		Обозначим последовательность Дика между первой $+1$ и найденной
		$-1$ как $x'$, а последовательность Дика после найденной $-1$ как
		$x''$. Любой начальный отрезок $x'$ имеет неотрицательную сумму,
		так как иначе минимальный префикс $x$ был бы меньше. Сумма по
		любому префиксу $x''$ также неотрицательна, а сумма по всему $x''$
		равна нулю.\bigskip

		Таким образом, для каждой последовательности Дика длины $2n$ мы
		сопоставили две последовательности Дика с длинами $2k$ и $2(n-k-1)$
		(два элемента из второй последовательности вычитаются, так как это
		первая $+1$ и найденная $-1$). При этом по парам таких
		последовательностей можно однозначно получить исходную, а значит
		это биекция. Тогда по правилу суммы выполняется равенство из условия.
	\end{proof}
\end{theorem}

\end{document}